\section{My role}

At 02:14 on 2026-03-31, after the second cancelled field day, I pushed commit \texttt{7eb7cc8} --- a fix for a double-scaling bug that had been quietly offsetting my GPS target estimates by more than 100 metres. I had been on the project for sixteen weeks by that point, and the commit was my 611th. What stopped me, staring at the diff, was not the bug itself but a smaller and more uncomfortable realisation: the file I was patching, \texttt{passive\_watch.py}, was not in my original remit. Neither were the 127 unit tests in \texttt{tests/unit/}, the \texttt{main.py} refactor from 1411 to 754 lines, the \texttt{simple\_simulator.py} MVP, or the D6 bibliography I would rebuild two days later. I had been allocated, at the wk12 kickoff, a narrow computer-vision role: train a YOLOv8n dummy detector and hand it to the flight team. \textbf{1a.}~Everything around that core --- the simulation framework, the progressive test ladder, the geofence logic, the blueprint system, the complaint letter to the course organiser, even Demetro's FDR speaking scripts --- I took on implicitly, one commit at a time, whenever I saw something between the team and a working demo. I initially framed this drift as conscientiousness (Hatton \& Smith, 1995). In hindsight I now see it was partly avoidance: building infrastructure alone was a way of feeling productive during the roughly eight weeks when three cancelled flight days left the team with no hardware to iterate against, and it protected me from the harder Western-team-engineering skill of slowing down so that a teammate could catch up. My Ukrainian training told me one person delivers; the project told me that is a failure mode dressed as a strength.

\textbf{1b.}~The contribution I am most proud of sits inside that same uncomfortable truth. It is not the retrained model at mAP50~=~0.995, nor the 820-plus commits, nor the D6 Goldmine climb from 74 to 85.2 across six scoring cycles. It is the fact that on the first cancelled flight day (2026-03-11), when the weather killed our only window, I did not go home --- I wrote \texttt{FIELD\_QUICK\_REF.md} for no-internet use, calibrated the focal length to 5.46 mm against a tape measure at 1~m, ran the checkerboard lens calibration, and closed ten commits before dark. That day crystallised my one-line thesis about this project, which I will defend: \emph{the simulation framework we built ourselves is the only reason this project produced any verifiable engineering output at all.} What I am proud of is not the output --- it is that I refused to let a missing drone become a missing report. The counterfactual is concrete: had I treated the field days as the unit of work, the team would have delivered zero flights \emph{and} zero verifiable data. Instead we delivered zero flights and a working state machine, a 2508-line simulator with GPS drift and camera shake, a 146-test pytest harness, and a web ground station a pilot could operate unaided. Had I acted on the ``build a simulator first'' instinct a week earlier than I did, we would have saved roughly five days of idle waiting --- a pattern I now recognise in myself: I trust the escape route too late.

\textbf{1c.}~What this experience has forced me to reframe, and what I will focus on differently in future, is the distinction between \emph{capability} and \emph{enablement}. I leave this project believing that engineering leadership is less about being the one who can fix the geofence sign-convention bug at midnight (reflection-in-action, Sch\"on, 1983), and more about the skill of making my fix legible enough that the teammate who owns the module would have found it without me. The evidence that convinced me is the \texttt{passive\_watch\_2.py} handoff to Robin (commit \texttt{a5b1ab7}): I built the \texttt{--robin} flag and the \texttt{cv\_mode} polling interface specifically so his state machine could consume my vision output over HTTP without touching my code --- and Robin then integrated for two days without asking me a single question. That was the first teamwork move I made on this project that scaled without my presence, and it happened in week 20, not week 12. In my next team project --- I start an internship with a robotics group in June --- I will schedule two hours per week of explicit pair-programming and handoff-document review, and the falsifiable signal I will watch for is whether a teammate can complete a handoff task independently within 48 hours of receiving it. If they cannot, my documentation has failed, not their attention. That is a sharper test than ``communicate better,'' and it is the test I should have been running on myself since wk12.
